\documentclass[11pt,a4paper]{article}

\usepackage[T1]{fontenc}
\usepackage[utf8]{inputenc}
\usepackage[french]{babel}
\usepackage[margin=2.5cm]{geometry}
\usepackage{booktabs}
\usepackage{siunitx}
\usepackage[table]{xcolor}
\usepackage{hyperref}
\usepackage{longtable}
\usepackage{pgfplots}
\pgfplotsset{
  compat=1.18,
  every axis legend/.append style={
    at={(0.5,0.98)}, anchor=north,
    legend columns=4,
    font=\scriptsize,
    fill=white, fill opacity=0.85, draw opacity=1, text opacity=1,
    draw=black!20,
    /tikz/every even column/.append style={column sep=6pt},
  },
}
\usepackage{caption}
\usepackage{enumitem}
\usepackage{amsmath}
\usepackage{amssymb}

\sisetup{detect-weight=true, detect-family=true}

\hypersetup{
  colorlinks=true,
  linkcolor=blue!50!black,
  urlcolor=blue!50!black,
  citecolor=blue!50!black,
  pdftitle={Rapport d'activite - Scan H(P) Al4C3},
  pdfauthor={Pr. John Akapulco}
}

\title{\vspace{-1.5cm}Rapport d'activité\\[2pt]
\large Diagramme de stabilité H(P) des polymorphes d'Al$_4$C$_3$ (0--100~GPa)}
\author{Pr. John Akapulco \\ \small \href{mailto:computchem86@gmail.com}{computchem86@gmail.com}}
\date{23 août 2026}

\begin{document}
\maketitle
\vspace{-0.5cm}

\begin{abstract}
\noindent
Ce rapport fait suite au rapport du 22~août~2026, qui comparait sept polymorphes candidats
d'Al$_4$C$_3$ à deux pressions seulement (0 et 50~GPa) et concluait sur un renversement de
stabilité entre R-3m (dominant à 0~GPa) et Pbam/Pnma (dominants à 50~GPa). La campagne décrite
ici étend la comparaison à une grille fine de onze pressions (0 à 100~GPa par pas de 10~GPa)
pour les sept polymorphes, soit 77 relaxations sous contrainte (\texttt{ISIF=8},
\texttt{PSTRESS}), chaînées séquentiellement par branche pour accélérer la convergence
électronique. Les 77 calculs ont convergé. Le diagramme de stabilité obtenu révèle trois
domaines successifs : \textbf{Pnma-VII} (0--26~GPa, quasi dégénéré avec R-3m en dessous de
20~GPa), \textbf{Pbam} (26--52~GPa) et \textbf{Pnma} (52--100~GPa), avec des pressions de
transition estimées à \textbf{26~GPa} et \textbf{52~GPa} (arrondies à l'unité de GPa ; voir
section~\ref{sec:precision} pour la justification). Le polymorphe \textbf{I-43d}
reste le moins stable sur toute la gamme mais converge rapidement vers les autres phases à
haute pression. Cette version ajoute les diagrammes $V(P)$ et $E(P)$ (avant celui de $H(P)$)
ainsi que le tableau numérique complet (énergie, volume, enthalpie par formule unitaire, au
millième d'eV/\AA$^3$) pour les 77 points du scan.

\medskip
\noindent\textbf{Mise à jour du 23~août 2026 (soir) :} le huitième polymorphe, \textbf{Fd-3m},
dont le scan PBE était encore en cours au moment de la rédaction initiale de ce rapport, a
depuis convergé sur l'intégralité de la grille 0--100~GPa (11/11 points), portant le total à
\textbf{88 points convergés sur 8 polymorphes}. Son ajout \emph{modifie} la séquence de
stabilité décrite plus bas : Fd-3m s'intercale entre Pnma-VII et Pbam et devient la phase la
plus stable sur une fenêtre intermédiaire, \textbf{5--32~GPa} environ, repoussant Pnma-VII à un
domaine plus étroit (0--5~GPa) et Pbam à 32--52~GPa ; la transition Pbam~$\to$~Pnma à 52~GPa
n'est en revanche pas affectée. Voir la section~\ref{sec:sequence}, révisée en conséquence.
\end{abstract}

\section{Contexte et objectifs}

Le rapport précédent (commit \texttt{97da5a3} et antérieurs) avait mis en évidence un
renversement d'ordre de stabilité entre R-3m et Pbam/Pnma entre 0 et 50~GPa, mais ne pouvait
ni localiser la pression de transition, ni exclure un comportement plus complexe entre ces deux
points. Les prochaines étapes qui y étaient formulées demandaient explicitement des points de
pression intermédiaires. La présente campagne y répond : un scan H(P) complet à onze pressions
(0, 10, 20, \dots, 100~GPa) a été lancé pour les sept polymorphes, afin de construire un
diagramme de stabilité continu et de localiser précisément les transitions de phase induites
par la pression.

\section{Méthodologie}

\subsection{Paramètres DFT}

Les paramètres reprennent ceux validés dans la campagne précédente, avec un maillage de
$\mathbf{k}$-points allégé pour ce balayage exploratoire :

\begin{center}
\begin{tabular}{ll}
\toprule
Paramètre & Valeur \\
\midrule
\texttt{PREC} & Accurate \\
\texttt{ENCUT} & 600~eV \\
\texttt{EDIFF} (convergence électronique) & $10^{-5}$~eV \\
\texttt{EDIFFG} (convergence ionique) & $-0.01$~eV/\AA\ (critère sur les forces) \\
\texttt{ISMEAR / SIGMA} & 0 / 0.05 \\
\texttt{KSPACING / KGAMMA} & 0.2~\AA$^{-1}$ / .TRUE. (maillage allégé pour le scan) \\
\texttt{IBRION / ISIF} & 2 / 8 (relaxation ions + volume, forme de maille fixée) \\
\texttt{NSW} & 200 \\
\texttt{PSTRESS} & 0 à 1000~kBar (0 à 100~GPa, pas de 100~kBar) \\
\texttt{KPAR} & 8 \\
\bottomrule
\end{tabular}
\end{center}

\subsection{Chaînage séquentiel par branche}

Pour chaque polymorphe, les onze points de pression sont soumis comme une seule tâche SLURM
(\texttt{run\_chain.sh}) qui exécute les onze relaxations \emph{séquentiellement dans la même
allocation}, en partant du point de départ le mieux relaxé disponible (0~GPa ou 50~GPa selon
la branche) puis en progressant par pas de 10~GPa de proche en proche. Chaque point redémarre
l'onde d'essai électronique (\texttt{WAVECAR}) et la densité de charge (\texttt{CHGCAR}) du
point de pression voisin déjà convergé, ce qui accélère nettement la convergence SCF par
rapport à un démarrage \texttt{ISTART=0} systématique. La table de dépendances
(\texttt{run\_manifest.txt}) trace pour chaque point la pression dont il dérive et, le cas
échéant, la structure candidate d'origine (\texttt{seed\_source}).

Les sept chaînes ont été soumises via \texttt{submit\_chain.sh} (jobs SLURM
58452--58458, un par polymorphe, partition \texttt{defq}) et se sont toutes terminées avec
succès. La chaîne Fd-3m, ajoutée séparément (\texttt{gen\_hp\_run\_fd3m.py}, chaîne ascendante
unique faute de \texttt{WAVECAR} convergé à toutes les pressions au moment du lancement), a
convergé de la même façon.

L'énergie totale $E$ (\texttt{"free~~energy~~~TOTEN"} dans l'\texttt{OUTCAR}) et le volume de
maille $V$ (\texttt{"volume~of~cell"}) sont extraits pour chaque point convergé en plus de
l'enthalpie $H$ ; les trois grandeurs sont rapportées par formule unitaire (division par
$N_{\text{atomes}}/7$).

\section{Résultats}

\subsection{État de convergence}

\begin{center}
\begin{tabular}{lc}
\toprule
Polymorphe & Points convergés (0--100~GPa) \\
\midrule
I-43d              & 11/11 \\
Pbam               & 11/11 \\
Pnma               & 11/11 \\
Pnma-III           & 11/11 \\
Pnma-VII           & 11/11 \\
R-3m               & 11/11 \\
anti-CaFe$_2$O$_4$ & 11/11 \\
Fd-3m              & 11/11 \\
\midrule
\textbf{Total} & \textbf{88/88} \\
\bottomrule
\end{tabular}
\end{center}

Convergence = critère de forces \texttt{EDIFFG=-0.01}~eV/\AA\ atteint (message \emph{"reached
required accuracy"} dans l'\texttt{OUTCAR}, terminaison confirmée par le bloc \emph{"General
timing"}).

\subsection{Diagramme de volume V(P)}

\begin{figure}[h]
\centering
\begin{tikzpicture}
\begin{axis}[
  width=\textwidth, height=8cm,
  xlabel={Pression (GPa)},
  ylabel={$V$/formule (\AA$^3$)},
  enlarge y limits={upper=0.22, lower=0.02},
  grid=major,
]
\addplot table[x=P_GPa,y=V_A3_per_fu] {../../hp_curves/hp_V_Pnma-VII.dat};
\addplot table[x=P_GPa,y=V_A3_per_fu] {../../hp_curves/hp_V_Pbam.dat};
\addplot table[x=P_GPa,y=V_A3_per_fu] {../../hp_curves/hp_V_Pnma.dat};
\addplot table[x=P_GPa,y=V_A3_per_fu] {../../hp_curves/hp_V_R-3m.dat};
\addplot table[x=P_GPa,y=V_A3_per_fu] {../../hp_curves/hp_V_Pnma-III.dat};
\addplot table[x=P_GPa,y=V_A3_per_fu] {../../hp_curves/hp_V_anti-CaFe2O4.dat};
\addplot table[x=P_GPa,y=V_A3_per_fu] {../../hp_curves/hp_V_I-43d.dat};
\addplot table[x=P_GPa,y=V_A3_per_fu] {../../hp_curves/hp_V_Fd-3m.dat};
\legend{Pnma-VII, Pbam, Pnma, R-3m, Pnma-III, anti-CaFe2O4, I-43d, Fd-3m}
\end{axis}
\end{tikzpicture}
\caption{Volume de la maille par formule unitaire (Al$_4$C$_3$) en fonction de la pression, pour
les huit polymorphes. Toutes les courbes décroissent de façon monotone et quasi parallèle : la
compressibilité est similaire d'un polymorphe à l'autre, sans changement de pente marqué au
voisinage des transitions de stabilité. Fd-3m (79.35~\AA$^3$/f.u.\ à 0~GPa) se situe dans le
haut du jeu, derrière R-3m et Pnma-VII, et suit la même loi de compression que les autres
branches.}
\label{fig:vp}
\end{figure}

\subsection{Diagramme d'énergie totale E(P)}

\begin{figure}[h]
\centering
\begin{tikzpicture}
\begin{axis}[
  width=\textwidth, height=8cm,
  xlabel={Pression (GPa)},
  ylabel={$E$/formule (eV)},
  enlarge y limits={upper=0.22, lower=0.02},
  grid=major,
]
\addplot table[x=P_GPa,y=E_eV_per_fu] {../../hp_curves/hp_E_Pnma-VII.dat};
\addplot table[x=P_GPa,y=E_eV_per_fu] {../../hp_curves/hp_E_Pbam.dat};
\addplot table[x=P_GPa,y=E_eV_per_fu] {../../hp_curves/hp_E_Pnma.dat};
\addplot table[x=P_GPa,y=E_eV_per_fu] {../../hp_curves/hp_E_R-3m.dat};
\addplot table[x=P_GPa,y=E_eV_per_fu] {../../hp_curves/hp_E_Pnma-III.dat};
\addplot table[x=P_GPa,y=E_eV_per_fu] {../../hp_curves/hp_E_anti-CaFe2O4.dat};
\addplot table[x=P_GPa,y=E_eV_per_fu] {../../hp_curves/hp_E_I-43d.dat};
\addplot table[x=P_GPa,y=E_eV_per_fu] {../../hp_curves/hp_E_Fd-3m.dat};
\legend{Pnma-VII, Pbam, Pnma, R-3m, Pnma-III, anti-CaFe2O4, I-43d, Fd-3m}
\end{axis}
\end{tikzpicture}
\caption{Énergie totale (sans le terme $PV$) par formule unitaire (Al$_4$C$_3$) en fonction de
la pression, pour les huit polymorphes. Contrairement à $H(P)$, l'ordre des courbes en $E(P)$
reste globalement celui de $P=0$ sur toute la gamme : c'est bien le terme $PV$ (piloté par les
écarts de volume, cf.\ figure~\ref{fig:vp}) qui provoque les croisements observés sur le
diagramme d'enthalpie.}
\label{fig:ep}
\end{figure}

\subsection{Diagramme de stabilité H(P)}

\begin{figure}[h]
\centering
\begin{tikzpicture}
\begin{axis}[
  width=\textwidth, height=8cm,
  xlabel={Pression (GPa)},
  ylabel={$H$/formule (eV)},
  enlarge y limits={upper=0.22, lower=0.02},
  grid=major,
]
\addplot table[x=P_GPa,y=H_eV_per_fu] {../../hp_curves/hp_Pnma-VII.dat};
\addplot table[x=P_GPa,y=H_eV_per_fu] {../../hp_curves/hp_Pbam.dat};
\addplot table[x=P_GPa,y=H_eV_per_fu] {../../hp_curves/hp_Pnma.dat};
\addplot table[x=P_GPa,y=H_eV_per_fu] {../../hp_curves/hp_R-3m.dat};
\addplot table[x=P_GPa,y=H_eV_per_fu] {../../hp_curves/hp_Pnma-III.dat};
\addplot table[x=P_GPa,y=H_eV_per_fu] {../../hp_curves/hp_anti-CaFe2O4.dat};
\addplot table[x=P_GPa,y=H_eV_per_fu] {../../hp_curves/hp_I-43d.dat};
\addplot table[x=P_GPa,y=H_eV_per_fu] {../../hp_curves/hp_Fd-3m.dat};
\legend{Pnma-VII, Pbam, Pnma, R-3m, Pnma-III, anti-CaFe2O4, I-43d, Fd-3m}
\end{axis}
\end{tikzpicture}
\caption{Enthalpie $H = E + PV$ par formule unitaire (Al$_4$C$_3$) en fonction de la pression,
pour les huit polymorphes. L'enveloppe inférieure (phase la plus stable à chaque pression) passe
successivement par Pnma-VII (0--5~GPa), \textbf{Fd-3m} (5--32~GPa), Pbam (32--52~GPa) puis Pnma
(52--100~GPa) -- Fd-3m, initialement absent de ce diagramme, s'y révèle être la phase la plus
stable sur près d'un tiers de la gamme de pression étudiée (voir section~\ref{sec:sequence}).}
\label{fig:hp}
\end{figure}

\subsection{Valeurs numériques complètes (E, V, H)}

Le tableau~\ref{tab:hp_full} rassemble, pour les 88 points convergés (huit polymorphes, Fd-3m
inclus), l'énergie totale $E$, le
volume $V$ et l'enthalpie $H$ par formule unitaire, arrondis au millième (eV pour $E$ et $H$,
\AA$^3$ pour $V$).

% Auto-generated from hp_curves/results_hp.csv -- do not edit by hand.
% Regenerate with: python3 hp_curves/extract_hp.py, then the table-export
% snippet in the report's methodology notes.
\begin{longtable}{lrS[table-format=-2.3]S[table-format=2.3]S[table-format=-2.3]}
\caption{Énergie totale $E$, volume $V$ et enthalpie $H = E + PV$ par formule unitaire
(Al$_4$C$_3$, 7 atomes), pour les sept polymorphes et les onze pressions du scan
(0--100~GPa). Valeurs arrondies au millième d'eV (resp.\ d'\AA$^3$).}
\label{tab:hp_full} \\
\toprule
Polymorphe & {$P$ (GPa)} & {$E$ (eV)} & {$V$ (\AA$^3$)} & {$H$ (eV)} \\
\midrule
\endfirsthead
\multicolumn{5}{l}{\small \textit{(suite)}} \\
\toprule
Polymorphe & {$P$ (GPa)} & {$E$ (eV)} & {$V$ (\AA$^3$)} & {$H$ (eV)} \\
\midrule
\endhead
\bottomrule
\endfoot
Pnma-VII & 0 & -43.354 & 81.237 & -43.354 \\
Pnma-VII & 10 & -43.219 & 76.785 & -38.426 \\
Pnma-VII & 20 & -42.893 & 73.287 & -33.745 \\
Pnma-VII & 30 & -42.445 & 70.410 & -29.261 \\
Pnma-VII & 40 & -41.912 & 67.972 & -24.943 \\
Pnma-VII & 50 & -41.316 & 65.845 & -20.768 \\
Pnma-VII & 60 & -40.680 & 63.987 & -16.717 \\
Pnma-VII & 70 & -40.004 & 62.325 & -12.773 \\
Pnma-VII & 80 & -39.301 & 60.828 & -8.929 \\
Pnma-VII & 90 & -38.577 & 59.465 & -5.173 \\
Pnma-VII & 100 & -37.835 & 58.215 & -1.499 \\
\midrule
Pbam & 0 & -42.226 & 73.355 & -42.226 \\
Pbam & 10 & -42.112 & 69.580 & -37.769 \\
Pbam & 20 & -41.837 & 66.625 & -33.520 \\
Pbam & 30 & -41.456 & 64.180 & -29.439 \\
Pbam & 40 & -40.997 & 62.080 & -25.498 \\
Pbam & 50 & -40.489 & 60.265 & -21.681 \\
Pbam & 60 & -39.937 & 58.655 & -17.971 \\
Pbam & 70 & -39.356 & 57.225 & -14.354 \\
Pbam & 80 & -38.739 & 55.915 & -10.820 \\
Pbam & 90 & -38.112 & 54.735 & -7.366 \\
Pbam & 100 & -37.460 & 53.635 & -3.983 \\
\midrule
Pnma & 0 & -41.596 & 71.517 & -41.596 \\
Pnma & 10 & -41.479 & 67.642 & -37.257 \\
Pnma & 20 & -41.201 & 64.653 & -33.130 \\
Pnma & 30 & -40.820 & 62.212 & -29.171 \\
Pnma & 40 & -40.372 & 60.160 & -25.352 \\
Pnma & 50 & -39.876 & 58.388 & -21.654 \\
Pnma & 60 & -39.337 & 56.815 & -18.061 \\
Pnma & 70 & -38.771 & 55.422 & -14.557 \\
Pnma & 80 & -38.179 & 54.160 & -11.136 \\
Pnma & 90 & -37.567 & 53.010 & -7.790 \\
Pnma & 100 & -36.944 & 51.960 & -4.512 \\
\midrule
R-3m & 0 & -43.329 & 81.533 & -43.329 \\
R-3m & 10 & -43.195 & 77.073 & -38.384 \\
R-3m & 20 & -42.869 & 73.577 & -33.684 \\
R-3m & 30 & -42.420 & 70.697 & -29.183 \\
R-3m & 40 & -41.887 & 68.253 & -24.846 \\
R-3m & 50 & -41.293 & 66.137 & -20.654 \\
R-3m & 60 & -40.654 & 64.270 & -16.586 \\
R-3m & 70 & -39.978 & 62.607 & -12.625 \\
R-3m & 80 & -39.276 & 61.110 & -8.763 \\
R-3m & 90 & -38.551 & 59.747 & -4.990 \\
R-3m & 100 & -37.810 & 58.497 & -1.299 \\
\midrule
Pnma-III & 0 & -42.360 & 74.380 & -42.360 \\
Pnma-III & 10 & -42.239 & 70.368 & -37.847 \\
Pnma-III & 20 & -41.947 & 67.228 & -33.554 \\
Pnma-III & 30 & -41.547 & 64.665 & -29.439 \\
Pnma-III & 40 & -41.072 & 62.495 & -25.470 \\
Pnma-III & 50 & -40.540 & 60.595 & -21.630 \\
Pnma-III & 60 & -39.975 & 58.947 & -17.900 \\
Pnma-III & 70 & -39.372 & 57.465 & -14.266 \\
Pnma-III & 80 & -38.745 & 56.127 & -10.719 \\
Pnma-III & 90 & -38.099 & 54.913 & -7.253 \\
Pnma-III & 100 & -37.436 & 53.797 & -3.859 \\
\midrule
anti-CaFe$_2$O$_4$ & 0 & -41.502 & 72.005 & -41.502 \\
anti-CaFe$_2$O$_4$ & 10 & -41.387 & 68.213 & -37.130 \\
anti-CaFe$_2$O$_4$ & 20 & -41.108 & 65.213 & -32.967 \\
anti-CaFe$_2$O$_4$ & 30 & -40.722 & 62.735 & -28.975 \\
anti-CaFe$_2$O$_4$ & 40 & -40.260 & 60.617 & -25.126 \\
anti-CaFe$_2$O$_4$ & 50 & -39.746 & 58.785 & -21.401 \\
anti-CaFe$_2$O$_4$ & 60 & -39.192 & 57.167 & -17.783 \\
anti-CaFe$_2$O$_4$ & 70 & -38.608 & 55.730 & -14.259 \\
anti-CaFe$_2$O$_4$ & 80 & -38.007 & 54.450 & -10.819 \\
anti-CaFe$_2$O$_4$ & 90 & -37.381 & 53.270 & -7.457 \\
anti-CaFe$_2$O$_4$ & 100 & -36.744 & 52.197 & -4.164 \\
\midrule
I-43d & 0 & -40.232 & 70.138 & -40.232 \\
I-43d & 10 & -40.096 & 65.405 & -36.014 \\
I-43d & 20 & -39.817 & 62.415 & -32.026 \\
I-43d & 30 & -39.449 & 60.050 & -28.204 \\
I-43d & 40 & -39.017 & 58.075 & -24.518 \\
I-43d & 50 & -38.539 & 56.370 & -20.947 \\
I-43d & 60 & -38.029 & 54.883 & -17.476 \\
I-43d & 70 & -37.488 & 53.553 & -14.091 \\
I-43d & 80 & -36.926 & 52.355 & -10.784 \\
I-43d & 90 & -36.345 & 51.263 & -7.549 \\
I-43d & 100 & -35.750 & 50.260 & -4.380 \\
\midrule
Fd-3m & 0 & -43.293 & 79.353 & -43.293 \\
Fd-3m & 10 & -43.166 & 75.109 & -38.478 \\
Fd-3m & 20 & -42.853 & 71.764 & -33.895 \\
Fd-3m & 30 & -42.422 & 69.005 & -29.501 \\
Fd-3m & 40 & -41.908 & 66.656 & -25.267 \\
Fd-3m & 50 & -41.335 & 64.619 & -21.169 \\
Fd-3m & 60 & -40.716 & 62.820 & -17.191 \\
Fd-3m & 70 & -40.063 & 61.214 & -13.319 \\
Fd-3m & 80 & -39.383 & 59.763 & -9.542 \\
Fd-3m & 90 & -38.682 & 58.443 & -5.852 \\
Fd-3m & 100 & -37.963 & 57.233 & -2.241 \\

\end{longtable}


\subsection{Séquence de phases et pressions de transition}
\label{sec:sequence}

\textbf{Révisé le 23~août 2026 (soir) suite à la convergence de Fd-3m.} Le classement à sept
polymorphes (sans Fd-3m) donnait une séquence à trois domaines, Pnma-VII~$\to$~Pbam~$\to$~Pnma
avec des transitions à 26 et 52~GPa. L'ajout de Fd-3m révèle qu'il s'intercale entre Pnma-VII et
Pbam et devient lui-même la phase la plus stable sur une large fenêtre intermédiaire : la
séquence correcte à huit polymorphes comporte donc \emph{quatre} domaines.

\begin{center}
\begin{tabular}{lcc}
\toprule
Domaine de pression & Phase la plus stable & Remarque \\
\midrule
0--5~GPa    & Pnma-VII & quasi dégénéré avec R-3m ($\Delta H < 0.025$~eV/f.u.) \\
5--32~GPa   & transition Pnma-VII~$\to$~Fd-3m & croisement interpolé à \textbf{5~GPa} \\
            & \textbf{Fd-3m} & le plus stable de 10 à 30~GPa sur la grille \\
32--52~GPa  & transition Fd-3m~$\to$~Pbam & croisement interpolé à \textbf{32~GPa} \\
            & Pbam & Pnma-III très proche ($\Delta H = 0.052$~eV/f.u. à 50~GPa) \\
50--60~GPa  & transition Pbam~$\to$~Pnma & croisement interpolé à \textbf{52~GPa} (inchangé) \\
60--100~GPa & Pnma & écart croissant avec les autres phases \\
\bottomrule
\end{tabular}
\end{center}

Les deux nouvelles pressions de transition sont obtenues par la même interpolation linéaire de
$H(P)$ entre les points de grille encadrant le croisement (0/10~GPa et 30/40~GPa
respectivement), ce qui donne brut 5.4~GPa et 32.1~GPa. La transition Pbam~$\to$~Pnma à 52~GPa,
elle, n'implique pas Fd-3m (non compétitif dans cette gamme, voir plus bas) et reste inchangée
par rapport à la version précédente de ce rapport. Toutes les pressions sont arrondies à
l'unité de GPa (voir section~\ref{sec:precision} pour la justification, étendue aux deux
nouvelles transitions).

I-43d demeure le polymorphe le moins stable sur toute la gamme étudiée
($\Delta H = +3.122$~eV/f.u. à 0~GPa par rapport à Pnma-VII), mais son écart avec la phase la
plus stable se réduit fortement sous pression ($\Delta H = +0.132$~eV/f.u. à 100~GPa par rapport
à Pnma), ce qui en fait le second candidat le plus proche du bas du classement au-delà de
80~GPa. Fd-3m, à l'inverse, quitte le classement des phases compétitives après 32~GPa : son
écart à l'enveloppe inférieure croît de façon monotone jusqu'à $\Delta H = +2.271$~eV/f.u. à
100~GPa par rapport à Pnma. Il termine la gamme en position intermédiaire, plus stable que
R-3m et Pnma-VII (qui héritent des rangs les plus défavorables à haute pression, inversant leur
avantage initial à basse pression) mais nettement moins stable qu'I-43d, qui rejoint au
contraire le groupe de tête au-delà de 80~GPa (cf.\ ci-dessus).

\section{Discussion}

\subsection{Confirmation et affinement du renversement observé précédemment}

Le rapport du 22~août signalait un renversement entre R-3m (dominant à 0~GPa) et Pbam/Pnma
(dominants à 50~GPa), établi sur seulement deux points de pression. Le scan complet confirme ce
renversement mais en révèle la structure fine : R-3m n'est en réalité jamais strictement le
plus stable des sept candidats sur la grille étudiée -- Pnma-VII le devance de peu dès 0~GPa
($\Delta H = 0.024$~eV/f.u.) et l'écart se creuse progressivement jusqu'à 20~GPa
($\Delta H = 0.060$~eV/f.u.). La
« domination » de R-3m relevée dans le rapport précédent provenait du jeu de données alors
partiel à 0~GPa, dans lequel Pnma-VII n'avait pas encore convergé. Le renversement
Pbam~$\to$~Pnma vers 50~GPa, en revanche, se confirme et se précise : il s'agit en fait de trois
transitions successives une fois Fd-3m intégré (Pnma-VII~$\to$~Fd-3m à 5~GPa, puis
Fd-3m~$\to$~Pbam à 32~GPa, puis Pbam~$\to$~Pnma à 52~GPa) et non d'une transition unique
R-3m~$\to$~Pbam/Pnma, ni même de la transition directe Pnma-VII~$\to$~Pbam à 26~GPa initialement
rapportée avant la convergence de Fd-3m.

\subsection{Précision des pressions de transition}
\label{sec:precision}

Les écarts d'enthalpie entre phases concurrentes autour des deux transitions restent modestes
(quelques dizaines à environ une centaine de meV/f.u.), en particulier entre Pbam et Pnma-III
sur toute la plage 30--90~GPa ($\Delta H < 0.113$~eV/f.u.). Le maillage de
$\mathbf{k}$-points allégé (\texttt{KSPACING=0.2}) utilisé pour ce scan exploratoire, plus
grossier que celui de la campagne 0/50~GPa (\texttt{KSPACING=0.03}), introduit une incertitude
numérique dont l'ordre de grandeur n'est pas mesuré directement ici, mais qui est vraisemblablement
supérieure au meV/atome.

Pour quantifier l'effet d'une telle incertitude sur la pression de transition elle-même, on
utilise la relation thermodynamique $dH/dP = V$ : au voisinage d'un croisement $P_t$ entre deux
phases, $\Delta H(P) = H_1(P) - H_2(P)$ a pour pente $d(\Delta H)/dP = \Delta V(P_t)$, le volume
relatif des deux phases à la transition. Une incertitude $\delta H$ sur chaque enthalpie se
propage donc en
\[
\delta P_t \approx \frac{\sqrt{2}\,\delta H}{|\Delta V(P_t)|},
\]
(facteur $\sqrt2$ pour des erreurs indépendantes sur les deux phases). Plus les deux phases en
compétition ont des volumes proches, plus la transition est sensible à l'incertitude sur $H$.

Appliqué aux trois transitions de ce scan (huit polymorphes), avec $\delta H = 1$~meV/atome
(borne optimiste) :

\begin{center}
\begin{tabular}{lccc}
\toprule
Transition & $\Delta V(P_t)$ (\AA$^3$/atome) & Pente $d(\Delta H)/dP$ & $\delta P_t$ \\
\midrule
Pnma-VII~$\to$~Fd-3m (5~GPa) & 0.253 & 1.58~meV/(atome$\cdot$GPa) & $\approx$0.63~GPa \\
Fd-3m~$\to$~Pbam (32~GPa)    & 0.682 & 4.26~meV/(atome$\cdot$GPa) & $\approx$0.24~GPa \\
Pbam~$\to$~Pnma (52~GPa)     & 0.267 & 1.67~meV/(atome$\cdot$GPa) & $\approx$0.85~GPa \\
\bottomrule
\end{tabular}
\end{center}

Ces trois valeurs sont des \emph{bornes optimistes} (elles supposent une précision d'enthalpie de
1~meV/atome, meilleure que ce que garantit \texttt{KSPACING=0.2}). Dans les trois cas, une
décimale sur la pression de transition n'est pas justifiée : \textbf{5~GPa}, \textbf{32~GPa} et
\textbf{52~GPa} (valeurs arrondies à l'unité de GPa) sont donc les seules pressions rapportées
dans ce document. La transition Pnma-VII~$\to$~Fd-3m est la moins bien déterminée des trois --
son $\Delta V$ est le plus petit et elle se situe dans la même région (0--10~GPa) où Pnma-VII et
R-3m sont déjà quasi dégénérés, ce qui en fait la transition la plus sensible à une reconvergence
au maillage fin. La quasi-dégénérescence Pbam/Pnma-III relevée ci-dessus subsiste par ailleurs sur
toute la plage 30--90~GPa ($\Delta H < 0.113$~eV/f.u., $\Delta V$ sous 0.03~\AA$^3$/atome au-delà
de 70~GPa), rendant tout ordre de stabilité entre elles non significatif à ce niveau de maillage.
Une confirmation avec le maillage fin de la campagne précédente (\texttt{KSPACING=0.03}) serait
nécessaire avant de considérer ces pressions comme définitives, en particulier pour la première
transition.

\section{Prochaines étapes}

\begin{enumerate}[label=(\roman*)]
  \item Reconverger les points de pression proches des trois transitions identifiées
  (0--10~GPa, 30--40~GPa et 50--60~GPa) avec le maillage fin (\texttt{KSPACING=0.03}) de la
  campagne précédente, pour confirmer les pressions de transition à 5, 32 et 52~GPa (voir
  section~\ref{sec:precision} pour l'analyse de précision qui justifie cet arrondi) --
  priorité à la première (0--10~GPa), la moins bien déterminée.
  \item Resserrer la grille de pression autour de ces trois transitions (pas de 2--5~GPa) pour
  réduire l'incertitude sur leur localisation, en particulier autour de 5~GPa où Fd-3m,
  Pnma-VII et R-3m sont proches.
  \item Calculer les phonons des quatre phases dominantes (Pnma-VII, Fd-3m, Pbam, Pnma) à leurs
  pressions de stabilité respectives, pour vérifier l'absence de modes imaginaires
  (stabilité dynamique) de part et d'autre des transitions.
  \item Exploiter les \texttt{DOSCAR} déjà produits à chaque point du scan pour suivre
  l'évolution du caractère électronique (métallique/isolant) le long des quatre branches.
\end{enumerate}

\appendix
\section{Références du dépôt}

Commit de lancement de la campagne : \texttt{97da5a3} (« Converge Pnma/anti-CaFe2O4/Pnma-VII
0GPa relaxations, launch 7-phase HP scan (0-100 GPa) »). Fd-3m ajouté au commit
\texttt{6d0b9bd} (« Complete Fd-3m PBE and r2SCAN H(P) scans »), amorcé séparément
(\texttt{gen\_hp\_run\_fd3m.py}). Données brutes des 88 relaxations dans
\texttt{hp\_curves/Al4C3\_<polymorphe>/P<pression>/}, scripts de soumission et d'extraction dans
\texttt{hp\_curves/} (\texttt{gen\_hp\_runs.py}, \texttt{submit\_chain.sh},
\texttt{extract\_hp.py}), résultats agrégés dans \texttt{hp\_curves/results\_hp.csv} et
\texttt{hp\_curves/hp\_<polymorphe>.dat} (H), \texttt{hp\_curves/hp\_E\_<polymorphe>.dat} (E),
\texttt{hp\_curves/hp\_V\_<polymorphe>.dat} (V). Tableau complet (section~3.5) généré depuis
\texttt{results\_hp.csv} dans \texttt{table\_hp\_full.tex} /
\texttt{table\_hp\_full\_body.tex}. Rapport précédent (campagne 0/50~GPa) :
\texttt{rapport\_activite\_Al4C3.tex}.

\end{document}
